\chapter*{Preface}
\addcontentsline{toc}{chapter}{Preface}

This book is written for a reader who can differentiate, integrate, and follow an algebraic calculation, but who has not yet been handed the usual stack of prerequisites for quantum field theory.  The aim is not to hide the subject's difficulty.  The aim is to spend that difficulty carefully.

Quantum field theory is often introduced after years of linear algebra, differential equations, special relativity, quantum mechanics, and group theory.  There is a good reason for that route: each tool earns its place.  Yet the core moves of the subject are not mysterious.  They are variations of functions, sums that become integrals, oscillators that become fields, symmetries that give conserved quantities, and local changes of phase that demand gauge fields.  Each of these begins in calculus.

The chapters therefore build the machinery as it is needed.  A vector is first a list of numbers with a length.  A matrix is first a rule for mixing lists.  A delta function is first a limit of narrow peaks tested under an integral.  A path integral is first an ordinary Gaussian integral repeated many times.  A gauge field is first the price paid when a symmetry is allowed to vary from point to point.  The advanced words are not avoided, but they are made to carry their weight.

The book is not a survey.  It is a course.  Equations are derived, not merely displayed.  When a step is formal, the text says so and shows the finite-dimensional calculation that motivates it.  When an argument depends on a physical assumption, that assumption is named.  Worked examples and exercises are included so that the reader can test each new idea with a pencil rather than by nodding at prose.

No short book can replace the long apprenticeship that makes a quantum field theorist.  Still, a self-contained first road through the subject can be honest.  That is the standard used here.

